% Problemas del capítulo 2. Generado por tools/problemas/gen.py; no editar a mano.
\section{Problemas}\label{sec:problemas-02}

Graduados: \(\star\) directo, \(\star\star\) para pensar, \(\star\star\star\) desafío. Las respuestas están en el apéndice~\ref{ap:respuestas}; conviene intentarlos antes de mirarlas.

\begin{problema}[\({\star}\)]\label{pb:02-1}
Al triángulo \(A=(0,0)\), \(B=(4,0)\), \(C=(0,3)\) aplicale primero la traslación de vector \(v=(2,-1)\) y después el giro de \(90^\circ\) alrededor del origen. Repetí en el otro orden. ¿Dónde termina \(C\) en cada caso?
\end{problema}

\begin{problema}[\({\star}{\star}\)]\label{pb:02-2}
Girás \(90^\circ\) alrededor de \(A=(0,0)\) y después \(-90^\circ\) alrededor de \(B=(4,0)\). Mostrá que el resultado es una traslación y encontrá su vector.
\end{problema}

\begin{problema}[\({\star}{\star}\)]\label{pb:02-3}
Escribí la matriz \(3\times3\) en coordenadas homogéneas del giro de \(30^\circ\) alrededor de \(c=(10,10)\).
\end{problema}

\begin{problema}[\({\star}{\star}\)]\label{pb:02-4}
Caminar y girar: arrancás en \((0,0)\) mirando al norte. Caminás \(10\)~m, girás \(90^\circ\) a la derecha, caminás \(10\)~m, girás \(45^\circ\) a la izquierda y caminás \(10\sqrt2\)~m. ¿Dónde terminás y hacia dónde mirás?
\end{problema}

\begin{problema}[\({\star}{\star}{\star}\)]\label{pb:02-5}
Escribí el punto \((3,4)\) en coordenadas polares. Dibujá las dos líneas coordenadas que pasan por él (la de \(r\) constante y la de \(\theta\) constante) y explicá por qué se cortan en ángulo recto.
\end{problema}
